\section{Introduction}

A self-replicating probe does not explore a galaxy alone. It settles a star, builds
copies of itself, and sends them on to further stars, so a single launch grows into a
front that sweeps outward and fills the reachable field. How long that takes is the
exploration-timescale question, and the sharpest treatment of it for self-replicators is
Nicholson and Forgan's slingshot study \cite{nicholson-forgan-2013}: probes that steal
energy from the motion of the stars they pass (a gravitational slingshot, following the
single-probe analysis of Forgan, Papadogiannakis and Kitching
\cite{forgan-papadog-kitching-2013}) can explore roughly two orders of magnitude faster
than probes under their own power.

That result rests on one convenient assumption, which its authors name as a limitation
and leave for future work: every probe is granted perfect, instantaneous knowledge of
which stars are already settled. No probe ever flies to a star another probe has just
claimed, because it always knows. In a real galaxy no such oracle exists. News that a
star has been settled can travel no faster than light, so a probe choosing its next
target is acting on a picture of the field that is, for distant stars, years or
centuries out of date.

This paper adds that light-speed limit to the model and measures what it costs. The
question is not merely academic bookkeeping: it decides whether the celebrated slingshot
speed-up survives contact with the finite speed of information, or whether a swarm racing
outward on stale maps wastes much of the advantage colliding with itself. The answer,
built from a controlled paired-ensemble experiment, is that the tax is real but sharply
selective - it is set in scale by the probe's speed and in bite by how far each hop
reaches.
